\section*{Exercise 1.1 -- 1.4}
\label{p1:sec:main}

\subsection*{Ex. 1.1 (5 pts) --- Magic number of a binary executable}
\label{p1:sec:magic}

\textbf{Role.} A magic number is a short, fixed byte pattern that a file carries so
that its \emph{format} can be recognised from its contents rather than from its name.
When the kernel services an \texttt{execve()}, it does not trust the file extension: it
reads the first handful of bytes and matches them against the formats it knows, and that
match is what selects the right loader. If nothing matches, the call is simply refused with \texttt{ENOEXEC}
instead of the machine jumping into arbitrary bytes. It therefore does two jobs at once:
format dispatch and a cheap integrity/sanity check.

\textbf{Where it is stored.} In the very first bytes of the file, at offset~0, as the
opening field of the file header, which is exactly what makes the check cheap, since
the kernel already has to read that first block anyway. e.g.: an ELF binary starts
with \texttt{0x7F 'E' 'L' 'F'} (the first four bytes of \texttt{e\_ident} in the ELF
header); a shell or Python script starts with \texttt{\#!} (\texttt{0x23 0x21}) followed
by the interpreter path; a Windows PE image starts with \texttt{'M' 'Z'}; a Java class
file starts with \texttt{0xCAFEBABE}.

\subsection*{Ex. 1.2 (5 pts) --- Segments of an address space}
\label{p1:sec:segments}

A process address space is built from the following regions, conventionally laid out from
low addresses to high:

\begin{itemize}[nosep,leftmargin=1.4em,label=$\bullet$]
  \item \textbf{Text (code) segment}: the machine instructions of the program. Mapped
        read-only and executable, so it can be shared by every process running the same
        binary and cannot be scribbled over by a stray pointer. Fixed size.
  \item \textbf{Initialized data segment (\texttt{.data})}: global and static variables
        whose initial value is not zero. Their values are stored in the executable file
        and copied in by the loader. Read/write, fixed size.
  \item \textbf{Uninitialized data segment (BSS)}: global and static variables that
        start out zero. Only their \emph{size} is recorded in the file, not their
        contents; the loader zero-fills the region, so the BSS costs nothing on disk.
        Read/write, fixed size.
  \item \textbf{Heap}: memory obtained at run time by \texttt{malloc}/\texttt{new}.
        It sits above the BSS and grows upward, toward higher addresses, as the
        allocator extends it (\texttt{brk}/\texttt{sbrk}, or \texttt{mmap} for large
        requests).
  \item \textbf{Memory-mapped region}: shared libraries, files mapped with
        \texttt{mmap}, and large allocations. It occupies the middle ground between heap
        and stack, and is how dynamically linked code is brought into the space.
  \item \textbf{Stack}: one per thread; holds activation records: parameters, local
        variables, saved registers and return addresses. It is placed at the top of the
        space and grows downward, toward lower addresses, one frame per call.
\end{itemize}

\subsection*{Ex. 1.3 (5 pts) --- Von Neumann machines and machines today}
\label{p1:sec:vonneumann}

A von Neumann machine is the \emph{stored-program} design: a CPU (control unit plus ALU
and registers), a single main memory that holds instructions \emph{and} data
indistinguishably, and I/O, all joined by a shared bus. Its defining idea is that the
program is itself just data sitting in memory, so the machine is general purpose and it operates by
repeating the "fetch->decode->execute" cycle, with a program counter naming the next
instruction. The price of the single shared memory and bus is that instruction fetch and
data access compete for the same path, the well-known \emph{von Neumann bottleneck}.

Today's machines are still von Neumann at the level the programmer observe, like x86, ARM
etc, all present one flat address space holding both code and data, a program
counter, and sequential instruction semantics. 

\subsection*{Ex. 1.4 (5 pts) --- ``Segmentation fault (core dumped)''}
\label{p1:sec:segfault}

\textbf{What the message means.} The program touched memory it was not entitled to
touch. The MMU checks every address as it translates it; if the address is unmapped,
or mapped with the wrong permission (writing a read-only page, executing a
non-executable one), it raises a hardware fault. The kernel finds no legitimate
reason for it, converts it into \texttt{SIGSEGV}, and delivers it to the process,
whose default action is to die and write a core dump. The line you see is the shell
afterwards reporting how its child died.

\textbf{Possible reasons.}
\begin{itemize}[nosep,leftmargin=1.4em,label=$\bullet$]
  \item Dereferencing a null pointer, typically an unchecked return from
        \texttt{malloc} or \texttt{fopen}.
  \item Dereferencing an uninitialized or dangling pointer, including use-after-free.
  \item Running past the end of an array far enough to leave the mapping.
  \item Writing through a \texttt{char *} that points at a string literal, which
        lives on a read-only page.
  \item Unbounded recursion overflowing the stack.
\end{itemize}

\textbf{What ``segmentation'' means here.} The word is historical. On segmented
machines the address space was divided into segments, each with a base, a limit and
access rights, and the hardware trapped any reference falling outside them: a
\emph{segmentation violation}. Mainstream systems are paged rather than segmented
now, but the name stuck and still means an access outside a valid, permitted region.

\textbf{What ``core'' means here.} ``Core'' is magnetic-core memory, used to be the ferrite-ring
main memory of machines of the 1950s. Then it became a synonym for main memory and
outlived the technology. A core dump is an image of the process's memory and register
state at the instant it faulted, written to a file so a debugger can reconstruct where
the program died.

